By the SCM moment-reduction lemma, under every \((\mu,\mathcal M)\in\mathfrak M_{m,\epsilon}\), the cluster counts are iid with categorical law \(b\), so \(Z\sim\operatorname{Multinomial}(n,b)\), and \(b\in\mathcal B_{m,\epsilon}\).  The size bound in the finite-complexity inference theorem, applied at the true \(b\), gives
\[
 \Pr_{\mathcal M}\{b\notin\mathcal R_{n,\alpha}(Z)\}
 =\Pr_{\mathcal M}\{P_n(b;Z)\leq\alpha\}\leq\alpha.
\]
On the complementary event, \(b\) is one of the candidates in \(\mathcal R_{n,\alpha}(Z)\cap\mathcal B_{m,\epsilon}\).  The defining union for \(\mathcal C_{n,\alpha}(Z)\) therefore contains its entire sharp fiber:
\[
 [L_m(b),U_m(b)]\subseteq\mathcal C_{n,\alpha}(Z).
\]
Taking probabilities proves the \(1-\alpha\) guarantee.  The argument uses no law-specific constant, hence is uniform over \(\mathfrak M_{m,\epsilon}\).  It invokes neither relative interior nor positive width and thus includes boundary laws with singleton fibers.